% ---------------------------------------------------
\chapter{Post-Link ELF Patch Tool}
\label{chap:postlink-elf-patch}
% ---------------------------------------------------

\section{Overview}
This chapter specifies the post-link utility that finalizes FLCR metadata
\emph{after} the linker has assigned physical addresses. The tool reads the
linked ELF, derives function \emph{sizes} from the symbol table, and computes
per-replica \emph{hashes} over code bytes. It then writes these values into the
\textbf{Function Size Table (FST)} and \textbf{Function Hash Table (FHT)}
corresponding to the replica addresses enumerated in the
\textbf{Function Address Table (FAT)}.

\paragraph{Inputs.}
A fully linked ELF binary containing:
\begin{itemize}
  \item \texttt{FAT} symbol (Function Address Table) — array of function start virtual addresses;
  \item \texttt{FST} symbol (Function Size Table) — writable array (u32) to be populated;
  \item \texttt{FHT} symbol (Function Hash Table) — writable array (u32) to be populated;
  \item \texttt{.symtab} (preferred) or \texttt{.dynsym} (fallback).
\end{itemize}

\paragraph{Outputs.}
The same ELF binary with \texttt{FST} and \texttt{FHT} entries updated in place.

\paragraph{Conventions.}
All examples below assume 32-bit virtual addresses for table entries and a
little-endian 32-bit word hash (\texttt{xor32}) with natural modulo-$2^{32}$
overflow and \emph{no tail padding}.

% ---------------------------------------------------
\section{Algorithms}

\begin{algorithm}[H]
\caption{\texttt{func\_size\_by\_vaddr(vaddr, Symtab)} — Linear scan of \texttt{.symtab}}
\label{alg:func-size-by-vaddr}
\begin{algorithmic}[1]
\State \textbf{Input:} \texttt{vaddr} (u32), \texttt{Symtab} (array of \texttt{Elf\_Sym})
\State \textbf{Output:} \texttt{size} (u32)
\For{\texttt{i} $\gets$ 0 \textbf{to} \texttt{Symtab.size() - 1}}
  \State \texttt{sym} $\gets$ \texttt{Symtab[i]}
  \If{\texttt{sym.st\_type} $==$ \texttt{STT\_FUNC}}
    \If{\texttt{vaddr} $==$ \texttt{sym.st\_value}}  \Comment{exact start match}
      \State \Return \texttt{sym.st\_size}
    \EndIf
  \EndIf
\EndFor
\State \Return \texttt{0} \Comment{no exact match found}
\end{algorithmic}
\end{algorithm}

\begin{algorithm}[H]
\caption{\texttt{populate\_fst(ELF)} — Fill Function Size Table (FST) from FAT via \texttt{.symtab}}
\label{alg:populate-fst}
\begin{algorithmic}[1]
\State \textbf{Input:} \texttt{ELF} (binary with tables \texttt{FAT}, \texttt{FST})
\State \textbf{Output:} \texttt{FST} updated in \texttt{ELF}
\State \texttt{Symtab} $\gets$ \textproc{find\_symtab}(\texttt{ELF}) \Comment{prefer \texttt{.symtab}, fallback \texttt{.dynsym}}
\State \texttt{FAT}    $\gets$ \textproc{find\_symbol}(\texttt{ELF}, \texttt{"function\_Address\_Table"})
\State \texttt{FST}    $\gets$ \textproc{find\_symbol}(\texttt{ELF}, \texttt{"function\_Size\_Table"})
\State \texttt{fat\_bytes} $\gets$ \textproc{symbol\_size}(\texttt{FAT})
\State \texttt{N} $\gets$ \texttt{fat\_bytes} $/ 4$ \Comment{u32 entries}
\State \texttt{func\_vaddr} $\gets$ \textproc{read\_u32\_array}(\texttt{ELF}, \texttt{FAT}) \Comment{array of \texttt{uint32\_t} vaddrs}
\For{\texttt{i} $\gets$ 0 \textbf{to} \texttt{N - 1}}
  \State \texttt{a} $\gets$ \texttt{func\_vaddr[i]} \Comment{function start vaddr}
  \State \texttt{n} $\gets$ \textproc{func\_size\_by\_vaddr}(\texttt{a}, \texttt{Symtab}) \Comment{linear scan of \texttt{.symtab}}
  \State \textproc{write\_u32\_at}(\texttt{ELF}, \texttt{FST}, \texttt{i}, \texttt{n})
\EndFor
\State \textproc{save}(\texttt{ELF})
\end{algorithmic}
\end{algorithm}

\begin{algorithm}[H]
\caption{\texttt{xor32(buf)} — 32-bit little-endian XOR, no tail padding}
\label{alg:xor32-no-pad}
\begin{algorithmic}[1]
\State \textbf{Input:} \texttt{buf} (byte array)
\State \textbf{Output:} \texttt{h} (u32)
\State \texttt{x} $\gets$ \texttt{0}
\For{\texttt{j} $\gets$ 0 \textbf{to} \textproc{len}(\texttt{buf}) $- 1$}
  \State \texttt{lane} $\gets$ \texttt{(j mod 4)} \Comment{byte position within 32-bit word}
  \State \texttt{x} $\gets$ \texttt{x} $\oplus$ \big(\texttt{buf[j]} $\ll$ \texttt{(8 $\times$ lane)}\big)
\EndFor
\State \Return \texttt{x} \Comment{natural overflow modulo $2^{32}$}
\end{algorithmic}
\end{algorithm}

\begin{algorithm}[H]
\caption{\texttt{populate\_fht(ELF)} — Fill Function Hash Table (FHT) using FAT (vaddrs) and FST (sizes)}
\label{alg:populate-fht}
\begin{algorithmic}[1]
\State \textbf{Input:} \texttt{ELF} (binary with tables \texttt{FAT}, \texttt{FST}, \texttt{FHT})
\State \textbf{Output:} \texttt{FHT} updated in \texttt{ELF}
\State \texttt{FAT} $\gets$ \textproc{find\_symbol}(\texttt{ELF}, \texttt{"function\_Address\_Table"})
\State \texttt{FST} $\gets$ \textproc{find\_symbol}(\texttt{ELF}, \texttt{"function\_Size\_Table"})
\State \texttt{FHT} $\gets$ \textproc{find\_symbol}(\texttt{ELF}, \texttt{"function\_Hash\_Table"})
\State \texttt{fat\_bytes} $\gets$ \textproc{symbol\_size}(\texttt{FAT}); \;\; \texttt{N} $\gets$ \texttt{fat\_bytes} $/ 4$
\State \texttt{func\_vaddr} $\gets$ \textproc{read\_u32\_array}(\texttt{ELF}, \texttt{FAT})
\State \texttt{sizes}      $\gets$ \textproc{read\_u32\_array}(\texttt{ELF}, \texttt{FST})
\State \textbf{assert} \textproc{symbol\_size}(\texttt{FST}) $==$ \texttt{N} $\times 4$
\State \textbf{assert} \textproc{symbol\_size}(\texttt{FHT}) $\ge$ \texttt{N} $\times 4$
\For{\texttt{i} $\gets$ 0 \textbf{to} \texttt{N - 1}}
  \State \texttt{a} $\gets$ \texttt{func\_vaddr[i]}; \;\; \texttt{n} $\gets$ \texttt{sizes[i]}
  \State \texttt{buf} $\gets$ \textproc{read\_bytes\_at\_vaddr}(\texttt{ELF}, \texttt{a}, \texttt{n})
  \State \texttt{h} $\gets$ \textproc{xor32}(\texttt{buf}) \Comment{no tail padding; natural modulo-$2^{32}$ overflow}
  \State \textproc{write\_u32\_at}(\texttt{ELF}, \texttt{FHT}, \texttt{i}, \texttt{h})
\EndFor
\State \textproc{save}(\texttt{ELF})
\end{algorithmic}
\end{algorithm}

% ---------------------------------------------------
\section{Notes and Assumptions}
\begin{itemize}
  \item \textbf{Symbol source:} prefer \texttt{.symtab}; fall back to \texttt{.dynsym} if necessary.
  \item \textbf{Exact-match policy:} sizes are derived only for symbols with
        \emph{exact} start address matches (\texttt{st\_value} equals FAT entry).
  \item \textbf{Word size / endianness:} examples assume 32-bit entries (u32) and little-endian hash lanes.
        Adapt \texttt{read\_u32\_array} / \texttt{write\_u32\_at} for 64-bit targets.
  \item \textbf{Section placement:} tables are updated in-place at the symbol’s
        section/offset as reported by the ELF symbol table.
\end{itemize}

% ---------------------------------------------------
\section{Summary}
The post-link tool completes the FLCR metadata by (i) deriving function sizes
from the linked symbol table and writing them to the FST, and (ii) hashing the
linked code bytes to populate the FHT. Together with the linker-assigned
addresses in the FAT, these tables enable runtime validation and selection of
replicas during \texttt{code\_init()}.